% GENERATED FILE. Edit canonical lessons, then run npm run generate:books.
% canonical-source-hash: fnv1a64:1e2f5b22f94dd91e
% canonical-lessons: FR-C175-embrasser, FR-C175-saluer, FR-C175-remercier, FR-C175-inviter, FR-C175-garer

\chapter{To kiss, To greet, To thank, To invite, To park}
\label{ch:fr-to-kiss-to-greet-to-thank-to-invite-to-park}
\hlchaptermodality{175}

\begin{chapteropening}
\textbf{By the end of this chapter:} \emph{I can say to kiss, to greet, to thank, to invite and to park.}

Complete the last lesson of chapter 175: I can say to kiss, to greet, to thank, to invite and to park.
\end{chapteropening}

\section[embrasser]{embrasser — to kiss}

\label{lesson:FR-C175-embrasser}



\begin{quote}
Before the new one: say the French for a meeting, then the French for a contract.
\end{quote}

\subsection*{You'll want to know: embrasser}
\textbf{embrasser} — \textquotedblleft{}to kiss\textquotedblright{}.

\begin{cousinweb}
\textbf{embrasser} — built from \emph{em-}, \textquotedblleft{}into\textquotedblright{}, and \emph{bras}, \textquotedblleft{}arm\textquotedblright{}: to take into the arms, now to kiss.
\end{cousinweb}

\begin{grammarlens}[title={What you've built}]
One more word for this chapter.
\end{grammarlens}

\subsection*{Your turn}
\begin{itemize}
\raggedright
  \item \emph{Say it:} \emph{embrasser}
  \item \emph{Say it:} \emph{embrasser}, once more
  \item \emph{Say it:} say \emph{embrasser}
  \item \emph{Recall:} say the French for a meeting, then the French for a contract, then say \emph{embrasser} again
\end{itemize}

\begin{tcolorbox}[breakable,colback=teal!4,colframe=teal!35!black,title={Before you move on}]
What does embrasser mean? (\textquotedblleft{}To kiss\textquotedblright{}.) Say it once more.
\end{tcolorbox}
\section[saluer]{saluer — to greet}

\label{lesson:FR-C175-saluer}



\begin{quote}
Before the new one: say the French for a contract, then the French for to kiss.
\end{quote}

\subsection*{You'll want to know: saluer}
\textbf{saluer} — \textquotedblleft{}to greet\textquotedblright{}.

\begin{cousinweb}
\textbf{saluer} $\leftarrow$ Latin \emph{salutare}, \textquotedblleft{}to greet, to wish health\textquotedblright{}, from \emph{salus}, \textquotedblleft{}health\textquotedblright{} — inherited. (EtymDB 2.1)
\end{cousinweb}

\begin{grammarlens}[title={What you've built}]
One more word for this chapter.
\end{grammarlens}

\subsection*{Your turn}
\begin{itemize}
\raggedright
  \item \emph{Say it:} \emph{saluer}
  \item \emph{Say it:} \emph{saluer}, once more
  \item \emph{Say it:} say \emph{saluer}
  \item \emph{Recall:} say the French for a contract, then the French for to kiss, then say \emph{saluer} again
\end{itemize}

\begin{tcolorbox}[breakable,colback=teal!4,colframe=teal!35!black,title={Before you move on}]
What does saluer mean? (\textquotedblleft{}To greet\textquotedblright{}.) Say it once more.
\end{tcolorbox}
\section[remercier]{remercier — to thank}

\label{lesson:FR-C175-remercier}



\begin{quote}
Before the new one: say the French for to kiss, then the French for to greet.
\end{quote}

\subsection*{You'll want to know: remercier}
\textbf{remercier} — \textquotedblleft{}to thank\textquotedblright{}.

\begin{cousinweb}
\textbf{remercier} — built from \emph{re-} and the old verb for \textquotedblleft{}to thank\textquotedblright{}, made from \emph{merci}, \textquotedblleft{}thank you\textquotedblright{}. (EtymDB 2.1 has its Middle French form.)
\end{cousinweb}

\begin{grammarlens}[title={What you've built}]
One more word for this chapter.
\end{grammarlens}

\subsection*{Your turn}
\begin{itemize}
\raggedright
  \item \emph{Say it:} \emph{remercier}
  \item \emph{Say it:} \emph{remercier}, once more
  \item \emph{Say it:} say \emph{remercier}
  \item \emph{Recall:} say the French for to kiss, then the French for to greet, then say \emph{remercier} again
\end{itemize}

\begin{tcolorbox}[breakable,colback=teal!4,colframe=teal!35!black,title={Before you move on}]
What does remercier mean? (\textquotedblleft{}To thank\textquotedblright{}.) Say it once more.
\end{tcolorbox}
\section[inviter]{inviter — to invite}

\label{lesson:FR-C175-inviter}



\begin{quote}
Before the new one: say the French for to greet, then the French for to thank.
\end{quote}

\subsection*{You'll want to know: inviter}
\textbf{inviter} — \textquotedblleft{}to invite\textquotedblright{}.

\begin{cousinweb}
\textbf{inviter} $\leftarrow$ Latin \emph{invitare}, \textquotedblleft{}to invite\textquotedblright{} — a learned borrowing. (EtymDB 2.1)
\end{cousinweb}

\begin{grammarlens}[title={What you've built}]
One more word for this chapter.
\end{grammarlens}

\subsection*{Your turn}
\begin{itemize}
\raggedright
  \item \emph{Say it:} \emph{inviter}
  \item \emph{Say it:} \emph{inviter}, once more
  \item \emph{Say it:} say \emph{inviter}
  \item \emph{Recall:} say the French for to greet, then the French for to thank, then say \emph{inviter} again
\end{itemize}

\begin{tcolorbox}[breakable,colback=teal!4,colframe=teal!35!black,title={Before you move on}]
What does inviter mean? (\textquotedblleft{}To invite\textquotedblright{}.) Say it once more.
\end{tcolorbox}
\section[garer]{garer — to park}

\label{lesson:FR-C175-garer}



\begin{quote}
Before the new one: say the French for to thank, then the French for to invite.
\end{quote}

\subsection*{You'll want to know: garer}
\textbf{garer} — \textquotedblleft{}to park\textquotedblright{}.

\begin{cousinweb}
\textbf{garer} $\leftarrow$ Old French \emph{garer}, \emph{warer}, \textquotedblleft{}to shelter, to guard\textquotedblright{} (EtymDB 2.1), from a Germanic verb meaning \textquotedblleft{}to watch over\textquotedblright{} — borrowed: to park is to put in shelter.
\end{cousinweb}

\begin{grammarlens}[title={What you've built}]
One more word for this chapter.
\end{grammarlens}

\subsection*{Your turn}
\begin{itemize}
\raggedright
  \item \emph{Say it:} \emph{garer}
  \item \emph{Say it:} \emph{garer}, once more
  \item \emph{Say it:} say \emph{garer}
  \item \emph{Recall:} say the French for to thank, then the French for to invite, then say \emph{garer} again
\end{itemize}

\begin{tcolorbox}[breakable,colback=teal!4,colframe=teal!35!black,title={Before you move on}]
What does garer mean? (\textquotedblleft{}To park\textquotedblright{}.) Say it once more.
\end{tcolorbox}
